\section{Matrix cones large for bounded traces}
\label{sec:tracial}

From this section onward, unless stated otherwise, $A$ is a separable
nuclear $C^*$-algebra satisfying \textup{(C)}.

The matrix cones in this section are approximately central in tracial
$2$-norm.  The construction uses a convex nuclear approximation, and its
estimates are independent of the number of maps in that approximation.
Throughout, $D\subseteq A$ is hereditary and $T_{\leq1}(D)$ denotes the
positive bounded traces on $D$ of norm at most one.  We extend these traces
normally to $D^{**}$.  On matrices we use the unnormalized extension
$\mathcal T=\Tr\otimes\tau$, and write
$\|x\|_{2,\tau}=\mathcal T(x^*x)^{1/2}$, also for rectangular matrices.
The traces need not be faithful.  All estimates below hold simultaneously
for every $\tau\in T_{\leq1}(D)$.

\begin{proposition}[Bounded-trace matrix cones]\label{prop:bounded-trace}
Fix $k\geq2$ and put
\[
 L=(20k)^2,\qquad C_k=1+4L.
\]
Let $F\subset D$ be a finite set of selfadjoint contractions, let
$\mathcal P\subseteq F$ consist of positive contractions, and let
$0<\delta<1$.  For every $\varepsilon>0$ there is a \cpc order-zero map
$V:M_k\to D$ such that, for every $\tau\in T_{\leq1}(D)$,
\begin{align}
 \sup_{\|x\|\leq1}\|[a,V(x)]\|_{2,\tau}
 &<\varepsilon &&(a\in F),\label{tr:central-target}\\
 \tau(aV(1))
 &\geq \frac\delta3\tau(a)-C_k\delta^{3/2}-\varepsilon
 &&(a\in\mathcal P).\label{tr:mass-target}
\end{align}
In particular, $k$ is fixed before the accuracy is chosen; no bound on the
number of nuclear approximation summands is required.
\end{proposition}

We first record the functional-calculus estimate used to regularize a
row whose operator norm need not be uniformly bounded.

\begin{lemma}[Lipschitz calculus in tracial $2$-norm]\label{tr:lipschitz}
Let $\mathcal M$ be a von Neumann algebra with a finite normal trace $T$,
let $H=H^*\in\mathcal M$, and let $B\in\mathcal M$.  If $f$ is real-valued
and Lipschitz on an interval containing the spectrum of $H$, then
\[
 \|[f(H),B]\|_{2,T}
 \leq\operatorname{Lip}(f)\|[H,B]\|_{2,T}.
\]
Consequently, if $W$ is a finite row, $a=a^*$ and $m=m^*$ are compatible
contractions, and
\[
 R=W(1+W^*W)^{-1/2},
\]
then $\|R\|\leq1$ and
\[
 \|aR-Rm\|_{2,T}\leq\|aW-Wm\|_{2,T}.
\]
The constants do not depend on the size of the matrices or on $\|W\|$.
\end{lemma}
\begin{proof}
For a finite spectral resolution $H=\sum_j\alpha_j e_j$, the blocks
$e_iBe_j$ are orthogonal in $L^2(\mathcal M,T)$.  Thus
\begin{align*}
 \|[f(H),B]\|_{2,T}^2
 &=\sum_{i,j}|f(\alpha_i)-f(\alpha_j)|^2\|e_iBe_j\|_{2,T}^2\\
 &\leq\operatorname{Lip}(f)^2
       \sum_{i,j}|\alpha_i-\alpha_j|^2\|e_iBe_j\|_{2,T}^2.
\end{align*}
Approximate an arbitrary $H$ in operator norm by finite spectral step
functions of $H$.  The corresponding functions $f(H)$ also converge in
norm.  Since $T(1)<\infty$, passage to $2$-norm proves the assertion.
This passage is made separately in each finite matrix algebra; the
resulting inequality has no factor involving its size or $T(1)$.

For the row assertion use
\[
 H=\begin{pmatrix}0&W\\W^*&0\end{pmatrix},\qquad
 B=\begin{pmatrix}a&0\\0&m\end{pmatrix},\qquad
 f(t)=\frac{t}{\sqrt{1+t^2}}.
\]
This function has Lipschitz constant one, and
$f(H)=\left(\begin{smallmatrix}0&R\\R^*&0\end{smallmatrix}\right)$.
The squares of the two commutator norms in the preceding inequality are
respectively $2\|aR-Rm\|_{2,T}^2$ and
$2\|aW-Wm\|_{2,T}^2$.
\end{proof}

We shall also use the following description of a finite-dimensional
commutant.  If $h\in D_+$ and
$\rho:E\to M(\her_D(h))$ is a unital homomorphism from a finite-dimensional
algebra, with unit $p_h$ in $D^{**}$, write
$E=\bigoplus_jM_{n_j}$ and let $v^{(j)}_{rs}$ be its represented matrix
units.  Then
\begin{equation}\label{tr:commutant-corners}
 \her_D(h)\cap\rho(E)'
 \cong\bigoplus_j v^{(j)}_{11}\her_D(h)v^{(j)}_{11},\qquad
 (b_j)_j\longmapsto\sum_{j,r}v^{(j)}_{r1}b_jv^{(j)}_{1r}.
\end{equation}
Indeed, commuting with the diagonal matrix units kills every unwanted
block, and commuting with the remaining matrix units identifies the
retained diagonal blocks by the displayed formula.  Each corner on the
right is hereditary in $D$.  Averaging over the compact unitary group
of $E$ gives a trace-preserving conditional expectation
\[
 \mathbb E_\rho:p_hD^{**}p_h\longrightarrow
                  p_hD^{**}p_h\cap\rho(E)'.
\]
It preserves $\her_D(h)$ because the represented unitaries are
multipliers.  It is contractive in tracial $2$-norm, by the triangle
inequality for the integral.  If $b$ is in its range, then
$\tau(b\mathbb E_\rho(x))=\tau(bx)$ whenever these expressions are defined.
These observations allow us to apply Lemma~\ref{lem:placement} in the
commutant, simultaneously to all its bounded positive traces.

\begin{proof}[Proof of Proposition~\ref{prop:bounded-trace}]
We will combine small matrix cones chosen in the commutants of the
nuclear approximation maps. Their supports are chosen to make the
model and overlap errors summable, after which a regularized row gives
a contraction. A second support compression makes this row bounded
below on the retained support. The polar isometry of the compressed
row then transports an exact matrix cone to $D$ while retaining the
tracial mass needed for the stated estimate.

Fix $F,\mathcal P,k,\delta$.  Our approximation parameter is
$0<u<1/4$.  The notation $o_u(1)$ below always means a quantity tending
to zero uniformly over $T_{\leq1}(D)$ with this fixed data.  It has no
dependence on the number of approximation summands.

\medskip\noindent
\textbf{The convex approximation and its weighted errors.}
The hereditary algebra $D$ is nuclear. The convex order-zero return
approximation is due to Hirshberg, Kirchberg, and White
\cite[Theorem~1.4]{HKW2012}; it is also part of the stronger
Brown--Carri\'on--White approximation theorem
\cite[Theorem~3.1]{BCW2019}. This supplies a finite-dimensional unital algebra
$E$, one \cpc map $\alpha:D\to E$, positive numbers
$\lambda_1,\ldots,\lambda_r$ with $\sum_l\lambda_l=1$, and \cpc
order-zero maps $\beta_l:E\to D$, such that
\[
 \Phi=\sum_{l=1}^r\lambda_l\beta_l\alpha
\]
is as close to the identity as prescribed on a given finite set.
The integer $r$ may depend on $u$ without restriction.  Put
\[
 h_l=\beta_l(1),\quad p_l=p_{h_l},\quad
 \beta_l(x)=h_l\rho_l(x),\quad
 \mu_l=\rho_l\alpha,\quad\theta_l=h_l\mu_l.
\]
Here $\rho_l$ is the supporting homomorphism of the order-zero map
\cite{WZ2009}; it acts by multipliers of $\her_D(h_l)$.
In particular $h_l$ commutes with $\mu_l(D)$.

By approximating $F$, its squares, and a sufficiently late positive
contractive approximate-unit element $e\in D$, we arrange
\begin{align}
 \left\|\sum_l\lambda_l(a-\theta_l(a))^2\right\|
 &\leq u^4 &&(a\in F),\label{tr:mean-square}\\
 \left\|a^{1/2}\left(1-\sum_l\lambda_lh_l\right)a^{1/2}\right\|
 &\leq u^4 &&(a\in\mathcal P).\label{tr:unit-error}
\end{align}
For completeness, the Schwarz inequality gives
\[
 0\leq\sum_l\lambda_l(a-\theta_l(a))^2
 \leq a^2-a\Phi(a)-\Phi(a)a+\Phi(a^2),
\]
whose norm is at most $3v$ if both relevant approximation errors are
at most $v$.  For the second estimate choose $e$ so that
$\|a^{1/2}(1-e)a^{1/2}\|\leq u^4/2$ for all targets, and approximate
$e$ within $u^4/2$.  Since $\theta_l(e)\leq h_l$, this proves
\eqref{tr:unit-error}.  All appearances of $1$ in this proof are allowed
in the bidual; no unit of $D$ is assumed.

Choose the following explicit continuous cutoff on $[0,1]$:
\[
 g(t)=\begin{cases}
 0,&t\leq1-2u,\\
 (t-1+2u)/u,&1-2u<t<1-u,\\
 1,&t\geq1-u.
 \end{cases}
\]
Set $g_l=g(h_l)$ and $z_l=1_{(1-2u,1]}(h_l)$.  Thus
$p_{g_l}\leq z_l\leq p_l$, and the scalar inequality
\begin{equation}\label{tr:g-loss}
 1-g_l^4\leq u^{-1}(1-h_l)
\end{equation}
holds by functional calculus.  With
\[
 D_l=\sum_{a\in F}p_l(a-\mu_l(a))^2p_l,\qquad
 B_u=2u^4+8u^2,
\]
we have, for every $a\in F$,
\begin{equation}\label{tr:model-square}
 \sum_l\lambda_l\|(a-\mu_l(a))z_l\|_{2,\tau}^2\leq B_u,
 \qquad
 \sum_l\lambda_l\tau(z_lD_l)\leq |F|B_u.
\end{equation}
Indeed, the contribution of $a-\theta_l(a)$ is bounded by $u^4$
using \eqref{tr:mean-square} and the trace inequality
$\tau(z_lxz_l)\leq\tau(x)$ for $x\geq0$.  On $z_l$,
$\|(\theta_l(a)-\mu_l(a))z_l\|\leq2u$.
Apply $\|x+y\|_2^2\leq2\|x\|_2^2+2\|y\|_2^2$ and
$\sum_l\lambda_l\tau(z_l)\leq1$.

Weighted Cauchy--Schwarz now gives, for $a\in\mathcal P$,
\begin{align}
 \sum_l\lambda_l\tau(\mu_l(a)g_l^4)
 &\geq\tau(a)-u^3-\sqrt{B_u},\label{tr:T1-lower}\\
 \sum_l\lambda_l\tau(\mu_l(a)z_l)
 &\leq\tau(a)+\sqrt{B_u}.\label{tr:T1-upper}
\end{align}
To verify the replacement of $\mu_l(a)$ by $a$, for example,
\[
 \left|\sum_l\lambda_l\tau((\mu_l(a)-a)g_l^4)\right|
 \leq
 \left(\sum_l\lambda_l\|(\mu_l(a)-a)z_l\|_2^2\right)^{1/2}
 \left(\sum_l\lambda_l\tau(g_l^8)\right)^{1/2}
 \leq\sqrt{B_u}.
\]
The same argument applies to $z_l$.  The lost $a$-mass in the first
inequality is at most $u^3$, by \eqref{tr:g-loss} and
\eqref{tr:unit-error}; in the second use $z_l\leq1$.

\medskip\noindent
\textbf{The first pencil: small supports and summable overlaps.}
For each $l$ put
\[
 C_l=\her_D(h_l)\cap\rho_l(E)',\qquad
 \mathbb E_l=\mathbb E_{\rho_l}.
\]
Construct maps $\psi_l:M_k\to C_l$ inductively, writing
$t_l=\psi_l(1)$. Once the earlier maps have been chosen, we penalize
both the model error and overlap with their weights by setting
\begin{align}
 S_l&=\sum_{i<l}t_i^2,\qquad
 K_l=\mathbb E_l\big(u^{-1}D_l+\delta^{-1/2}p_lS_lp_l\big),\nonumber\\
 X_l&=g_l(p_l-K_l)g_l,\qquad
 b_l=g_l(X_l)_+g_l.\label{tr:first-pencil}
\end{align}
These formulas define $b_l\in C_l$.  To check membership before taking
positive parts, sandwich each term of $D_l$ by $g_l$.  Terms involving
$a\in D$ then lie in $\her_D(h_l)$, and terms involving $\mu_l(a)$ lie
there by multiplier action.  The same is true of $g_lS_lg_l$.
Averaging preserves this hereditary algebra and commutes with
multiplication by $g_l$.  Hence $X_l$ belongs to it and to the
commutant, as does $b_l$.

The inequalities $X_l\leq p_l$ and $X_l\geq-(\|K_l\|+1)p_l$
imply $0\leq(X_l)_+\leq p_l$ by scalar functional calculus.  Therefore
\begin{equation}\label{tr:first-pencil-bounds}
 0\leq b_l\leq g_l^2\leq z_l,\qquad
 b_l\geq g_l^4-g_l^2K_lg_l^2.
\end{equation}
The second inequality follows from $(X_l)_+\geq X_l$.
Lemma~\ref{lem:support}, applied to the pencil in
\eqref{tr:first-pencil}, gives
\begin{equation}\label{tr:first-support}
 p_{b_l}K_lp_{b_l}\leq p_{b_l}.
\end{equation}

Use the tracial conclusion of Lemma~\ref{lem:placement} in $C_l$,
which has the form \eqref{tr:commutant-corners}, with
$q=\delta\lambda_l$ and positive contraction $b_l$.
It gives a \cpc order-zero map
$\psi_l:M_k\to\her_{C_l}(b_l)$ whose weight $t_l$ satisfies
\begin{equation}\label{tr:placement}
 d_\sigma(t_l)\leq\delta\lambda_l d_\sigma(b_l),\qquad
 d_\sigma((t_l-1/2)_+)
 \geq\frac{\delta\lambda_l}{3}d_\sigma((b_l-1/2)_+)
\end{equation}
for every bounded positive trace $\sigma$ on $C_l$.
In particular, on restricting $\tau$ to $C_l$,
\begin{equation}\label{tr:support-budget}
 \tau(p_{t_l})\leq\delta\lambda_l,\qquad
 \sum_l\tau(p_{t_l})\leq\delta,\qquad \tau(S_l)\leq\delta.
\end{equation}
Here ranks are the traces of the corresponding support projections in
the bidual.  As $p_{t_l}\leq p_{b_l}$, \eqref{tr:first-support} gives
\[
 \sum_l\tau(t_l^2K_l)\leq\sum_l\tau(t_l^2)\leq\delta.
\]
We may remove $\mathbb E_l$ inside these traces because $t_l^2\in C_l$.
Both summands defining $K_l$ are positive.  It follows that
\begin{align}
 \sum_l\|(a-\mu_l(a))t_l\|_{2,\tau}^2
 &\leq u\delta &&(a\in F),\label{tr:T2-model}\\
 \sum_l\tau(t_l^2S_l)&\leq\delta^{3/2}.
 \label{tr:T2-overlap}
\end{align}
These are unweighted sums: the individual support allowances
$\delta\lambda_l$ in \eqref{tr:support-budget} have absorbed the
convex coefficients. Their total is $\delta$, so increasing the number
of approximation summands does not enlarge either error bound.

We next obtain a lower bound for their high spectral cuts.  Let
\[
 s_l=1_{(1/2,1]}(t_l),\qquad e_l=1_{(1/2,1]}(b_l).
\]
Since $b_l$ is a contraction supported under $z_l$, scalar functional
calculus and \eqref{tr:first-pencil-bounds} give
\begin{equation}\label{tr:cut-lower}
 e_l\geq2b_l-z_l\geq2g_l^4-2g_l^2K_lg_l^2-z_l.
\end{equation}
For $a\in\mathcal P$, the functional
\[
 \sigma_{l,a}(c)=\tau(\mu_l(a)c)\qquad(c\in C_l)
\]
is a bounded positive trace.  Indeed $\mu_l(a)$ is a positive
contraction commuting with $C_l$, so positivity and the trace identity
follow from those of $\tau$.  Normality shows that
$d_{\sigma_{l,a}}(t_l)=\tau(\mu_l(a)p_{t_l})$, and gives the same
formula for every positive spectral cut.  Thus
\eqref{tr:placement} applies to this single weighted trace and yields
$\tau(\mu_l(a)s_l)\geq(\delta\lambda_l/3)\tau(\mu_l(a)e_l)$.
Using \eqref{tr:cut-lower} and $0\leq\mu_l(a)\leq1$, we obtain
\begin{align}
 \sum_l\tau(\mu_l(a)s_l)
 &\geq \frac\delta3\sum_l\lambda_l
               \tau(\mu_l(a)(2g_l^4-z_l))
       -\frac{2\delta}{3}\sum_l\lambda_l\tau(g_l^2K_lg_l^2).
 \label{tr:weighted-mass}
\end{align}
For the last term, remove the expectation against $g_l^4\in C_l$:
\begin{align*}
 \sum_l\lambda_l\tau(g_l^2K_lg_l^2)
 &=\sum_l\lambda_l\left(u^{-1}\tau(g_l^4D_l)
                   +\delta^{-1/2}\tau(g_l^4S_l)\right)\\
 &\leq u^{-1}|F|B_u+\sqrt\delta.
\end{align*}
In the inequality use $g_l^4\leq z_l$, \eqref{tr:model-square}, and
$\tau(S_l)\leq\delta$.  Combining this with
\eqref{tr:T1-lower}--\eqref{tr:T1-upper} proves
\begin{equation}\label{tr:T3}
 \sum_l\tau(\mu_l(a)s_l)
 \geq\frac\delta3\tau(a)-\delta^{3/2}-r_u
 \qquad(a\in\mathcal P),
\end{equation}
where one may take the explicit error
\[
 r_u=\frac\delta3\left(2u^3+3\sqrt{B_u}
                          +2|F|B_u/u\right)\longrightarrow0.
\]

\medskip\noindent
\textbf{Regularizing the combined row.}
Form matrices over $D$ and models in its bidual as follows:
\[
 W=(t_1,\ldots,t_r),\quad D_0=\operatorname{diag}(t_l^2),\quad
 m(a)=\operatorname{diag}(\mu_l(a)),\quad p=p_{D_0},
\]
\[
 R=W(1+W^*W)^{-1/2},\qquad Q_0=R^*R,\qquad
 G_0=D_0(1+D_0)^{-1}.
\]
Here $\mathcal T(p)\leq\delta$ by \eqref{tr:support-budget}.
Each $t_l$ commutes with $\mu_l(a)$, so \eqref{tr:T2-model} says
\[
 \|aW-Wm(a)\|_{2,\tau}\leq\sqrt{u\delta}=:e_u.
\]
The diagonal of $W^*W-D_0$ is zero, and its $(i,j)$ entry off the
diagonal is $t_it_j$.  Consequently
\begin{equation}\label{tr:off-diagonal-cost}
 \|W^*W-D_0\|_{2,\tau}^2
 =2\sum_l\sum_{i<l}\tau(t_l^2t_i^2)
 =2\sum_l\tau(t_l^2S_l)\leq2\delta^{3/2}.
\end{equation}
Lemma~\ref{tr:lipschitz} gives the first of the following estimates;
the second follows by expanding the commutator of $R^*R$:
\begin{align}
 \|aR-Rm(a)\|_{2,\tau}&\leq e_u,
 &\|[Q_0,m(a)]\|_{2,\tau}&\leq2e_u.
 \label{tr:regularized-central}
\end{align}
Finally, $Q_0=1-(1+W^*W)^{-1}$ and
$G_0=1-(1+D_0)^{-1}$.  The resolvent identity and the contractivity of
both resolvents show that
\begin{equation}\label{tr:regularized-overlap}
 \|Q_0-G_0\|_{2,\tau}^2\leq2\delta^{3/2}.
\end{equation}
This regularization has removed every occurrence of $\|W\|$ from the
estimates.

\medskip\noindent
\textbf{The second pencil in one matrix commutant.}
Write $B_0=\her_{M_r(D)}(D_0)$.  Both $W^*W$ and $D_0$ belong to
$B_0$: for the former, use its entries $t_it_j$ and an approximate
unit obtained by functional calculus of $D_0$.  Hence $Q_0,G_0\in B_0$.
Let
\[
 \pi:M_k\longrightarrow M(B_0)
\]
be the diagonal of the supporting homomorphisms of $\psi_l$, whose
unit in the bidual is $p$.  Its values commute with
$\operatorname{diag}(t_l)$ and with $m(a)$.  To see the latter assertion,
recover each supporting homomorphism by bounded strong limits from
$t_l$ and $\psi_l(M_k)$, all of which commute with $\mu_l(a)$.
The multiplier assertion follows because its products with
$\operatorname{diag}(t_l)$ belong to $M_r(D)$; the same observation
shows that $m(a)$ acts by multipliers on $B_0$ and commutes with $p$.

The row $R$ is contractive, but its Gram operator $Q_0$ need not
commute with $\pi(M_k)$, so $R\pi(\cdot)R^*$ need not have order zero.
We use the small overlap cost in \eqref{tr:regularized-overlap} to find
a $\pi(M_k)$-invariant support on which $R$ is bounded below. The polar
isometry of the compressed row will then transport the matrix action
while the estimates below control the lost tracial mass.

Let $u_k,v_k\in M_k$ be the clock and shift unitaries.  The map
\[
 \mathbb E_\pi(X)=\frac1{k^2}\sum_{i,j=0}^{k-1}
       \pi(u_k^iv_k^j)X\pi(u_k^iv_k^j)^*
\]
is the trace-preserving expectation onto the commutant of $\pi(M_k)$
in the corner with unit $p$.  The usual matrix-unit calculation shows
that the average kills the nonconstant matrix coefficients, and fixes
the commutant.  Since the identity conjugation is one of its $k^2$
positive summands,
\begin{equation}\label{tr:finite-average-bound}
 X\leq k^2\mathbb E_\pi(X)\qquad(X\geq0).
\end{equation}
Put
\[
 K=\mathbb E_\pi((Q_0-G_0)^2).
\]
Then $K\in B_0\cap\pi(M_k)'$, and
\begin{equation}\label{tr:second-cost}
 \mathcal T(K)\leq2\delta^{3/2},\qquad
 \|[m(a),K]\|_{2,\tau}\leq4e_u.
\end{equation}
For the commutator use $[m(a),G_0]=0$, $\|Q_0-G_0\|\leq1$,
\eqref{tr:regularized-central}, and the $2$-norm contractivity of the
expectation, which commutes with multiplication by $m(a)$.

Choose the fixed piecewise linear cutoff
\[
 g'(t)=\begin{cases}
 0,&t\leq1/10,\\
 10t-1,&1/10<t<1/5,\\
 1,&t\geq1/5,
 \end{cases}
 \qquad g'=g'(G_0),
\]
and, with $L=(20k)^2$, set
\begin{equation}\label{tr:second-pencil}
 X'=g'(p-LK)g',\qquad c=g'(X')_+g'.
\end{equation}
As for the first pencil, $c\in B_0\cap\pi(M_k)'$ and
\begin{equation}\label{tr:second-pencil-bounds}
 0\leq c\leq(g')^2\leq p,\qquad
 c\geq(g')^4-L(g')^2K(g')^2,
 \qquad p_c\leq1_{(1/10,\infty)}(G_0).
\end{equation}
Lemma~\ref{lem:support} gives $p_cKp_c\leq L^{-1}p_c$.
For $\xi\in p_c\mathcal H$, \eqref{tr:finite-average-bound} implies
\[
 \|(Q_0-G_0)\xi\|^2
 \leq k^2\langle K\xi,\xi\rangle
 \leq\frac1{400}\|\xi\|^2.
\]
On the same space $G_0$ has compression at least $1/10$.  Therefore
\begin{equation}\label{tr:polar-lower}
 c^2\geq cQ_0c\geq\frac1{20}c^2.
\end{equation}
Furthermore $m(a)$ commutes with $g'$.  Lemma~\ref{tr:lipschitz},
applied to the scalar function $t\mapsto t_+$ in
\eqref{tr:second-pencil}, yields the useful explicit bound
\begin{equation}\label{tr:c-commutator}
 \|[m(a),c]\|_{2,\tau}
 \leq L\|[m(a),K]\|_{2,\tau}\leq4Le_u.
\end{equation}

There remains a lower bound on the mass surviving this second pencil.
Set $s'=\operatorname{diag}(s_l)$.  On $s_l$ one has $t_l>1/2$,
so $t_l^2/(1+t_l^2)>1/5$, and hence $g's'=s'$.
The scalar inequality $c^2\geq2c-p$, followed by
\eqref{tr:second-pencil-bounds}, gives
\[
 s'c^2s'\geq2s'cs'-s'\geq s'-2Ls'Ks'.
\]
For $a\in\mathcal P$, $m(a)$ is a positive contraction commuting with
$s'$.  Decomposing the trace along $s'$ and $p-s'$ shows that
$\mathcal T(m(a)c^2)\geq\mathcal T(m(a)s'c^2s')$;
the omitted diagonal contribution is nonnegative.  Thus
\begin{align}
 \mathcal T(m(a)c^2)
 &\geq\mathcal T(m(a)s')-2L\mathcal T(K)\nonumber\\
 &\geq\frac\delta3\tau(a)-(1+4L)\delta^{3/2}-r_u.
 \label{tr:corrected-mass}
\end{align}
Here we used \eqref{tr:T3} and \eqref{tr:second-cost}.  The argument
only uses the commutation of $m(a)$ with $s'$; it does not require
$c$ to commute with $s'$.

\medskip\noindent
\textbf{Polar correction and exact order zero.}
Put $d=cQ_0c$ and let $U$ be the polar partial isometry of the row $Rc$
in the bidual.  Equation~\eqref{tr:polar-lower} implies
$p_d=p_c$, so $U^*U=p_c$.  We claim that the row
\[
 Y=Uc
\]
belongs to $M_{1,r}(D)$ and satisfies
\begin{equation}\label{tr:polar-intertwining}
 \sup_{\tau\in T_{\leq1}(D)}\|aY-Ym(a)\|_{2,\tau}=o_u(1)
 \qquad(a\in F).
\end{equation}
We give uniform estimates for this passage, since $r$ can grow with $u$.

For $\eta>0$ define the algebraic row
\[
 Y_\eta=Rc(d+\eta)^{-1/2}c\in M_{1,r}(D).
\]
The domination $c^2\leq20d$ implies a factorization
$c=d^{1/2}x$ with $\|x\|\leq\sqrt{20}$.  One elementary proof is to
define $x^*$ initially on the range of $d^{1/2}$ by
$x^*(d^{1/2}\xi)=c\xi$; the stated domination makes this well-defined
and bounded, and it extends by zero on the orthogonal complement.
Using $Rc=Ud^{1/2}$ now gives
\begin{equation}\label{tr:uniform-polar-norm}
 \|Y-Y_\eta\|
 \leq\sqrt{20}\sup_{t\geq0}
        \sqrt t\left(1-\sqrt{\frac{t}{t+\eta}}\right)
 \leq\sqrt{20\eta}.
\end{equation}
For the last scalar bound split at $t=\eta$: below it use $\sqrt t$,
and above it use
$1-\sqrt{t/(t+\eta)}\leq\eta/t$.
In particular the norm limit $Y$ lies in the rectangular algebra.
Both $Y$ and $Y_\eta$ have right support under $p$, so
\begin{equation}\label{tr:uniform-polar-two}
 \|Y-Y_\eta\|_{2,\tau}\leq\sqrt{20\eta\delta}.
\end{equation}
This uses $\mathcal T(p)\leq\delta$, rather than a bound on the trace
of the full matrix unit.

At fixed $\eta$ the intertwining error for $Y_\eta$ tends to zero
uniformly with $u$.  To spell this out, \eqref{tr:regularized-central}
and \eqref{tr:c-commutator} imply
\[
 \|[m(a),d]\|_{2,\tau}\leq(8L+2)e_u.
\]
On $[0,1]$ the function $(t+\eta)^{-1/2}$ has norm at most
$\eta^{-1/2}$ and Lipschitz constant at most
$\tfrac12\eta^{-3/2}$.  Expand $aY_\eta-Y_\eta m(a)$ through its four
factors and apply Lemma~\ref{tr:lipschitz}.  This gives
\[
 \|aY_\eta-Y_\eta m(a)\|_{2,\tau}
 \leq\left((1+8L)\eta^{-1/2}
                    +(4L+1)\eta^{-3/2}\right)e_u.
\]
Together with \eqref{tr:uniform-polar-two}, this bounds the error for
$Y$ by the displayed right-hand side plus $2\sqrt{20\eta\delta}$.
First choose $\eta$ small and then $u$ small.  This proves
\eqref{tr:polar-intertwining} without any restriction on $r$.

Define
\[
 V(x)=Uc\pi(x)cU^*=Y\pi(x)Y^*\qquad(x\in M_k).
\]
These elements lie in $D$: the rows $Y_\eta\pi(x)$ belong to
$M_{1,r}(D)$ by multiplier action, and converge in norm to
$Y\pi(x)$ by \eqref{tr:uniform-polar-norm}.
The map is completely positive and contractive, since
$V(1)=Uc^2U^*\leq1$.  It is exactly order zero.  Indeed $c$ commutes
with $\pi(M_k)$, $U^*U=p_c$, and for orthogonal positive $x,y\in M_k$,
\[
 V(x)V(y)=Uc^4\pi(xy)U^*=0.
\]
For a selfadjoint $a\in F$ set $\Delta_a=aY-Ym(a)$.  Since $m(a)$
commutes with $\pi(M_k)$ and $\|Y\|\leq1$,
\[
 [a,V(x)]=\Delta_a\pi(x)Y^*-Y\pi(x)\Delta_a^*,\qquad
 \|[a,V(x)]\|_{2,\tau}\leq2\|x\|\|\Delta_a\|_{2,\tau}.
\]
This proves \eqref{tr:central-target} when $u$ is sufficiently small.
Finally $Y^*Y=c^2$ and $\|Y\|_{2,\tau}\leq\sqrt\delta$.  Rectangular
traciality and Cauchy--Schwarz give
\[
 \left|\tau(aV(1))-\mathcal T(m(a)c^2)\right|
 =\left|\mathcal T(Y^*\Delta_a)\right|
 \leq\sqrt\delta\,\|\Delta_a\|_{2,\tau}=o_u(1).
\]
Combining this with \eqref{tr:corrected-mass}, and choosing $u$ to make
$r_u+\sqrt\delta\sup_\tau\|\Delta_a\|_{2,\tau}<\varepsilon$ for the
finitely many targets, proves \eqref{tr:mass-target}.  All choices were
made for the entire family $T_{\leq1}(D)$ at once, completing the proof.
\end{proof}
